\documentclass[11pt,a4paper]{article}

%====================================================
% PACKAGES
%====================================================
\usepackage[utf8]{inputenc}
\usepackage[T1]{fontenc}
\usepackage[french]{babel}
\usepackage[a4paper,margin=2cm]{geometry}
\usepackage{amsmath,amssymb}
\usepackage{lmodern}
\usepackage{fancyhdr}
\usepackage{tikz}
\usepackage{enumitem}

\pagestyle{fancy}
\fancyhf{}
\fancyhead[L]{Mathématiques}
\fancyhead[C]{Projection orthogonale}
\fancyhead[R]{\thepage}

\begin{document}

\begin{center}
\Huge \textbf{PROJECTION ORTHOGONALE}\\
\Large Cours complet + méthode + schémas
\end{center}

%====================================================
\section{Introduction et fondements géométriques}
%====================================================

\subsection{Motivation}

Dans de nombreux problèmes de mathématiques et de physique, on est amené à répondre à la question suivante :

\begin{center}
\textit{Quel est le point d’une droite ou d’un plan le plus proche d’un point donné ?}
\end{center}

\bigskip

Ce problème apparaît notamment dans :
\begin{itemize}
\item le calcul de distances,
\item la géométrie analytique,
\item la mécanique (forces, trajectoires),
\item l’algèbre linéaire (projections).
\end{itemize}

\bigskip

La réponse repose sur une notion fondamentale :

\begin{center}
\Large \textbf{la projection orthogonale.}
\end{center}

---

%====================================================
\subsection{Définition}
%====================================================

Soit un point $A$ et une droite $(d)$.

\bigskip

On appelle \textbf{projection orthogonale de $A$ sur $(d)$} le point $H$ tel que :

\[
H \in (d) \quad \text{et} \quad \vec{AH} \perp (d).
\]

Autrement dit, le segment $AH$ est perpendiculaire à la droite $(d)$.

---

%====================================================
\subsection{Représentation géométrique}
%====================================================

\begin{center}
\begin{tikzpicture}[scale=1.2]

% Droite
\draw[thick] (-3,0) -- (3,0) node[right] {$(d)$};

% Points
\filldraw (1.5,2) circle (2pt) node[above] {$A$};
\filldraw (1.5,0) circle (2pt) node[below] {$H$};

% Segment
\draw[thick] (1.5,2) -- (1.5,0);

% Angle droit
\draw (1.5,0) rectangle +(0.3,0.3);

\end{tikzpicture}
\end{center}

\bigskip

Le point $H$ est appelé le \textbf{projeté orthogonal} de $A$ sur $(d)$.

---

%====================================================
\subsection{Propriété fondamentale : distance minimale}
%====================================================

\textbf{Théorème.}  
Parmi tous les points de la droite $(d)$, le point $H$ est celui qui minimise la distance à $A$.

\bigskip

\textbf{Démonstration.}

Soit $M$ un point quelconque de la droite $(d)$.

\bigskip

Les points $A$, $H$ et $M$ forment un triangle rectangle en $H$, car $\vec{AH} \perp (d)$.

\bigskip

On applique le théorème de Pythagore :

\[
AM^2 = AH^2 + HM^2.
\]

Or $HM^2 \geq 0$, donc :

\[
AM^2 \geq AH^2.
\]

En prenant la racine :

\[
AM \geq AH.
\]

\bigskip

Ainsi, $H$ est le point de la droite $(d)$ le plus proche de $A$.

\hfill $\square$

---

%====================================================
\subsection{Interprétation analytique}
%====================================================

La condition de perpendicularité peut être traduite à l’aide du produit scalaire.

\bigskip

\textbf{Propriété.}  
Deux vecteurs sont orthogonaux si et seulement si leur produit scalaire est nul.

\bigskip

Ainsi, si $\vec{u}$ est un vecteur directeur de la droite $(d)$ :

\[
\vec{AH} \perp (d)
\quad \Longleftrightarrow \quad
\vec{AH} \cdot \vec{u} = 0.
\]

---

%====================================================
\subsection{Lecture géométrique du produit scalaire}
%====================================================

\begin{center}
\begin{tikzpicture}[scale=1.2]

\draw[->] (0,0) -- (3,0) node[right] {$\vec{u}$};
\draw[->] (0,0) -- (0,2.5) node[above] {$\vec{AH}$};

\node at (0.4,0.4) {$90^\circ$};

\end{tikzpicture}
\end{center}

\bigskip

Le produit scalaire permet donc de traduire une condition géométrique (orthogonalité)
en une équation algébrique.

---

%====================================================
\subsection{Vision globale}
%====================================================

La projection orthogonale peut être vue sous plusieurs angles :

\begin{center}
\begin{tabular}{|c|c|}
\hline
Vision géométrique & Vision analytique \\
\hline
perpendicularité & produit scalaire nul \\
\hline
distance minimale & optimisation \\
\hline
\end{tabular}
\end{center}

---

%====================================================
\subsection{Annonce de la suite}
%====================================================

Dans la suite du chapitre, on va :

\begin{itemize}
\item déterminer explicitement le point $H$ (projection sur une droite),
\item généraliser au cas d’un plan,
\item introduire une formulation vectorielle plus efficace (niveau prépa),
\item faire le lien avec le changement de base.
\end{itemize}

\bigskip

\begin{center}
\textbf{La projection orthogonale est une manière d’extraire une composante d’un vecteur.}
\end{center}
%====================================================
\section{Projection orthogonale sur une droite}
%====================================================

\subsection{Objectif}

On considère :

\begin{itemize}
\item un point $A$
\item une droite $(d)$
\end{itemize}

\bigskip

\begin{center}
\textbf{On cherche le point $H$ projection orthogonale de $A$ sur $(d)$}
\end{center}

---

%====================================================
\subsection{Données mathématiques}
%====================================================

Pour travailler, on suppose que la droite est donnée par :

\begin{itemize}
\item un point $B$ appartenant à $(d)$
\item un vecteur directeur $\vec{u}$
\end{itemize}

---

%====================================================
\subsection{Schéma fondamental}
%====================================================

\begin{center}
\begin{tikzpicture}[scale=1.3]

% droite
\draw[thick] (-3,0) -- (3,0) node[right] {$(d)$};

% points
\filldraw (1,2) circle (2pt) node[above] {$A$};
\filldraw (1,0) circle (2pt) node[below] {$H$};
\filldraw (-1,0) circle (2pt) node[below] {$B$};

% segments
\draw[thick] (1,2) -- (1,0);

% angle droit
\draw (1,0) rectangle +(0.25,0.25);

% vecteur directeur
\draw[->] (-1,0) -- (0,0) node[below] {$\vec{u}$};

\end{tikzpicture}
\end{center}

---

%====================================================
\subsection{Idée clé de la méthode}
%====================================================

Pour trouver $H$, on utilise deux informations :

\begin{enumerate}
\item $H$ appartient à la droite
\item $\vec{AH}$ est perpendiculaire à la droite
\end{enumerate}

---

%====================================================
\subsection{Étape 1 : paramétrisation de $H$}
%====================================================

Comme $H$ appartient à $(d)$ :

\[
\boxed{
H = B + t\vec{u}
}
\]

où $t$ est un réel.

---

%====================================================
\subsection{Étape 2 : condition de perpendicularité}
%====================================================

On traduit :

\[
\vec{AH} \perp (d)
\]

par :

\[
\boxed{
\vec{AH} \cdot \vec{u} = 0
}
\]

---

%====================================================
\subsection{Étape 3 : expression de $\vec{AH}$}
%====================================================

\[
\vec{AH} = \vec{OH} - \vec{OA}
\]

Or :

\[
\vec{OH} = \vec{OB} + t\vec{u}
\]

Donc :

\[
\vec{AH} = (\vec{OB} + t\vec{u}) - \vec{OA}
\]

\[
\boxed{
\vec{AH} = \vec{AB} + t\vec{u}
}
\]

---

%====================================================
\subsection{Étape 4 : équation en $t$}
%====================================================

On remplace dans la condition :

\[
(\vec{AB} + t\vec{u}) \cdot \vec{u} = 0
\]

Développement :

\[
\vec{AB} \cdot \vec{u} + t (\vec{u} \cdot \vec{u}) = 0
\]

---

\[
\boxed{
t = - \frac{\vec{AB} \cdot \vec{u}}{||\vec{u}||^2}
}
\]

---

%====================================================
\subsection{Étape 5 : expression finale de $H$}
%====================================================

\[
\boxed{
H = B - \frac{\vec{AB} \cdot \vec{u}}{||\vec{u}||^2} \vec{u}
}
\]

---

%====================================================
\subsection{Interprétation fondamentale}
%====================================================

Le vecteur $\vec{BH}$ est :

\[
\vec{BH} = - \frac{\vec{AB} \cdot \vec{u}}{||\vec{u}||^2} \vec{u}
\]

---

\begin{center}
\textbf{$\vec{BH}$ est la projection de $\vec{BA}$ sur $\vec{u}$}
\end{center}

---

%====================================================
\subsection{Schéma d’interprétation}
%====================================================

\begin{center}
\begin{tikzpicture}[scale=1.2]

\draw[->] (0,0) -- (4,0) node[right] {$\vec{u}$};

\draw[->] (0,0) -- (2,2) node[right] {$\vec{BA}$};

\draw[dashed] (2,2) -- (2,0);

\node at (2,-0.3) {$\vec{BH}$};

\end{tikzpicture}
\end{center}

---

%====================================================
\subsection{Exemple détaillé}
%====================================================

On considère :

\[
A(1,2), \quad B(0,0), \quad \vec{u} = (2,1)
\]

---

\textbf{1. Calcul de $\vec{AB}$}

\[
\vec{AB} = (1,2)
\]

---

\textbf{2. Produit scalaire}

\[
\vec{AB} \cdot \vec{u} = 1 \cdot 2 + 2 \cdot 1 = 4
\]

---

\textbf{3. Norme}

\[
||\vec{u}||^2 = 2^2 + 1^2 = 5
\]

---

\textbf{4. Coefficient}

\[
t = -\frac{4}{5}
\]

---

\textbf{5. Point projeté}

\[
H = B + t\vec{u}
= (0,0) + \left(-\frac{4}{5}\right)(2,1)
\]

\[
\boxed{
H = \left(-\frac{8}{5}, -\frac{4}{5}\right)
}
\]

---

%====================================================
\subsection{Conclusion}
%====================================================

\begin{itemize}
\item la méthode repose uniquement sur :
\begin{itemize}
\item paramétrisation
\item produit scalaire
\end{itemize}

\item aucune formule n’est à apprendre par cœur
\end{itemize}

\bigskip

\begin{center}
\textbf{Comprendre la méthode est bien plus important que mémoriser une formule.}
\end{center}

%====================================================
\section{Projection orthogonale sur un plan}
%====================================================

\subsection{Problématique}

On considère :

\begin{itemize}
\item un point $A$ de l’espace,
\item un plan $\mathcal{P}$.
\end{itemize}

\bigskip

\begin{center}
\textit{On cherche le point $H$ du plan $\mathcal{P}$ le plus proche de $A$.}
\end{center}

\bigskip

Ce point est appelé \textbf{projection orthogonale de $A$ sur le plan}.

---

%====================================================
\subsection{Définition}
%====================================================

Soit $\mathcal{P}$ un plan et $A$ un point.

\bigskip

Le point $H$ est la projection orthogonale de $A$ sur $\mathcal{P}$ si :

\[
H \in \mathcal{P}
\quad \text{et} \quad
\vec{AH} \perp \mathcal{P}.
\]

---

%====================================================
\subsection{Vecteur normal}
%====================================================

Un plan est caractérisé par un \textbf{vecteur normal} $\vec{n}$.

\bigskip

\textbf{Propriété fondamentale :}

\begin{center}
$\vec{AH} \perp \mathcal{P} \quad \Longleftrightarrow \quad \vec{AH}$ est colinéaire à $\vec{n}$.
\end{center}

\bigskip

Donc :

\[
\vec{AH} = t \vec{n}
\]

pour un certain réel $t$.

---

%====================================================
\subsection{Schéma géométrique}
%====================================================

\begin{center}
\begin{tikzpicture}[scale=1]

% plan
\filldraw[fill=gray!20] (-3,-1) -- (3,-1) -- (4,1) -- (-2,1) -- cycle;

% point A
\filldraw (0,3) circle (2pt) node[above] {$A$};

% point H
\filldraw (0,0) circle (2pt) node[below] {$H$};

% segment
\draw[thick] (0,3) -- (0,0);

% vecteur normal
\draw[->] (0,0) -- (0,1.5) node[right] {$\vec{n}$};

% angle droit
\draw (0,0) rectangle +(0.3,0.3);

\node at (2.5,0) {$\mathcal{P}$};

\end{tikzpicture}
\end{center}

---

%====================================================
\subsection{Méthode générale}
%====================================================

On suppose que le plan $\mathcal{P}$ a une équation :

\[
ax + by + cz + d = 0
\]

et donc un vecteur normal :

\[
\vec{n} = (a,b,c)
\]

---

\textbf{Étape 1 : paramétrer le point $H$}

\[
H = A + t\vec{n}
\]

---

\textbf{Étape 2 : imposer que $H$ appartient au plan}

On remplace les coordonnées de $H$ dans l’équation du plan.

---

\textbf{Étape 3 : résoudre en $t$}

---

\textbf{Étape 4 : déterminer $H$}

---

%====================================================
\subsection{Démonstration de la méthode}
%====================================================

On pose :

\[
A(x_A,y_A,z_A)
\]

\[
H = (x_A + ta,\; y_A + tb,\; z_A + tc)
\]

---

On impose que $H \in \mathcal{P}$ :

\[
a(x_A + ta) + b(y_A + tb) + c(z_A + tc) + d = 0
\]

---

Développement :

\[
ax_A + by_A + cz_A + d + t(a^2 + b^2 + c^2) = 0
\]

---

Donc :

\[
\boxed{
t = -\frac{ax_A + by_A + cz_A + d}{a^2 + b^2 + c^2}
}
\]

---

%====================================================
\subsection{Expression finale de la projection}
%====================================================

\[
\boxed{
H = A - \frac{ax_A + by_A + cz_A + d}{a^2 + b^2 + c^2} \vec{n}
}
\]

---

%====================================================
\subsection{Interprétation}
%====================================================

Le vecteur $\vec{AH}$ est la projection de $A$ sur la direction normale.

\bigskip

\begin{center}
\textbf{On se déplace depuis $A$ uniquement selon la direction perpendiculaire au plan.}
\end{center}

---

%====================================================
\subsection{Exemple détaillé}
%====================================================

On considère :

\[
A(1,2,3)
\]

\[
\mathcal{P} : x + y + z = 1
\]

---

\textbf{Vecteur normal :}

\[
\vec{n} = (1,1,1)
\]

---

\textbf{Calcul de $t$ :}

\[
t = -\frac{1+2+3-1}{1+1+1}
= -\frac{5}{3}
\]

---

\textbf{Point projeté :}

\[
H = (1,2,3) + \left(-\frac{5}{3}\right)(1,1,1)
\]

\[
\boxed{
H = \left(-\frac{2}{3}, \frac{1}{3}, \frac{4}{3}\right)
}
\]

---

%====================================================
\subsection{Distance point-plan}
%====================================================

La distance de $A$ au plan est :

\[
d = ||\vec{AH}||
\]

---

On obtient la formule classique :

\[
\boxed{
d = \frac{|ax_A + by_A + cz_A + d|}{\sqrt{a^2 + b^2 + c^2}}
}
\]

---

%====================================================
\subsection{Conclusion}
%====================================================

\begin{itemize}
\item la projection repose sur la direction normale
\item la méthode est analogue à celle d’une droite
\item la perpendicularité devient une colinéarité
\end{itemize}

\bigskip

\begin{center}
\textbf{Projection sur un plan = déplacement selon la normale}
\end{center}

%====================================================
\section{Distance minimale et interprétation de la projection}
%====================================================

\subsection{Problématique générale}

On considère un point $A$ et une droite $(d)$.

\bigskip

\begin{center}
\textit{On cherche à minimiser la distance entre $A$ et un point $M$ appartenant à $(d)$.}
\end{center}

\bigskip

Autrement dit, on cherche :

\[
\min_{M \in (d)} AM
\]

---

%====================================================
\subsection{Schéma d’interprétation}
%====================================================

\begin{center}
\begin{tikzpicture}[scale=1.2]

\draw[thick] (-3,0) -- (3,0);

\filldraw (1.5,2) circle (2pt) node[above] {$A$};
\filldraw (1.5,0) circle (2pt) node[below] {$H$};
\filldraw (-1,0) circle (2pt) node[below] {$M$};

\draw[thick] (1.5,2) -- (1.5,0);
\draw[dashed] (1.5,2) -- (-1,0);

\draw (1.5,0) rectangle +(0.3,0.3);

\node at (2.8,0) {$(d)$};

\end{tikzpicture}
\end{center}

---

%====================================================
\subsection{Résultat fondamental}
%====================================================

\textbf{Théorème :}

Le point $H$, projection orthogonale de $A$ sur $(d)$, est l’unique point qui minimise la distance à $A$.

---

%====================================================
\subsection{Démonstration (cas d’une droite)}
%====================================================

Soit $M$ un point quelconque de la droite $(d)$.

\bigskip

On considère les vecteurs :

\[
\vec{AM} = \vec{AH} + \vec{HM}
\]

---

On calcule :

\[
||\vec{AM}||^2 = (\vec{AH} + \vec{HM}) \cdot (\vec{AH} + \vec{HM})
\]

---

Développement :

\[
||\vec{AM}||^2 = ||\vec{AH}||^2 + ||\vec{HM}||^2 + 2\,\vec{AH} \cdot \vec{HM}
\]

---

Or :

\[
\vec{AH} \perp (d) \quad \text{et} \quad \vec{HM} \in (d)
\]

Donc :

\[
\vec{AH} \cdot \vec{HM} = 0
\]

---

On obtient :

\[
||\vec{AM}||^2 = ||\vec{AH}||^2 + ||\vec{HM}||^2
\]

---

Comme :

\[
||\vec{HM}||^2 \geq 0
\]

on a :

\[
||\vec{AM}||^2 \geq ||\vec{AH}||^2
\]

---

Donc :

\[
AM \geq AH
\]

---

\[
\boxed{
AH \text{ est la distance minimale}
}
\]

\hfill $\square$

---

%====================================================
\subsection{Interprétation analytique}
%====================================================

On peut reformuler ce problème comme un problème d’optimisation.

\bigskip

On paramètre la droite :

\[
M(t) = B + t\vec{u}
\]

---

On définit la fonction :

\[
f(t) = ||AM(t)||^2
\]

---

On cherche à minimiser $f(t)$.

---

\subsection{Calcul}

\[
f(t) = ||A - (B + t\vec{u})||^2
\]

---

On dérive :

\[
f'(t) = -2 \vec{u} \cdot (A - B - t\vec{u})
\]

---

On annule :

\[
f'(t) = 0 \quad \Longleftrightarrow \quad \vec{u} \cdot \vec{AH} = 0
\]

---

\[
\boxed{
\text{on retrouve la condition de projection orthogonale}
}
\]

---

%====================================================
\subsection{Cas du plan}
%====================================================

Le raisonnement est analogue.

\bigskip

On montre que la distance minimale est obtenue lorsque :

\[
\vec{AH} \parallel \vec{n}
\]

où $\vec{n}$ est un vecteur normal au plan.

---

%====================================================
\subsection{Vision unifiée}
%====================================================

\begin{center}
\begin{tabular}{|c|c|}
\hline
Géométrie & Analyse \\
\hline
perpendicularité & dérivée nulle \\
\hline
projection & minimum \\
\hline
\end{tabular}
\end{center}

---

%====================================================
\subsection{Interprétation profonde}
%====================================================

La projection orthogonale est la solution d’un problème d’optimisation :

\begin{center}
\textbf{minimiser une distance sous contrainte}
\end{center}

---

\subsection{Lien avec la suite}

Cette vision est essentielle pour comprendre :

\begin{itemize}
\item les projections en algèbre linéaire
\item les espaces vectoriels
\item les méthodes d’optimisation
\end{itemize}

\bigskip

\begin{center}
\textbf{La projection est une optimisation géométrique.}
\end{center}

%====================================================
\section{Vers la prépa : formulation vectorielle de la projection}
%====================================================

\subsection{Objectif}

Dans les parties précédentes, nous avons déterminé la projection d’un point sur une droite
à l’aide d’une méthode systématique (paramétrisation + produit scalaire).

\bigskip

Nous cherchons maintenant une formulation :

\begin{center}
\textbf{plus directe, plus rapide et plus générale.}
\end{center}

---

%====================================================
\subsection{Problème fondamental}
%====================================================

Soit :

\begin{itemize}
\item un vecteur $\vec{u} \neq 0$,
\item un vecteur $\vec{v}$.
\end{itemize}

\bigskip

On souhaite :

\begin{center}
\textit{projeter $\vec{v}$ sur la direction de $\vec{u}$.}
\end{center}

---

%====================================================
\subsection{Schéma}
%====================================================

\begin{center}
\begin{tikzpicture}[scale=1.2]

\draw[->] (0,0) -- (4,0) node[right] {$\vec{u}$};
\draw[->] (0,0) -- (2.5,2) node[right] {$\vec{v}$};

\draw[dashed] (2.5,2) -- (2.5,0);

\node at (2.5,-0.4) {projection};

\end{tikzpicture}
\end{center}

---

%====================================================
\subsection{Idée clé}
%====================================================

La projection de $\vec{v}$ sur $\vec{u}$ est un vecteur colinéaire à $\vec{u}$.

\bigskip

On cherche donc un réel $\lambda$ tel que :

\[
\text{proj}_{\vec{u}}(\vec{v}) = \lambda \vec{u}.
\]

---

%====================================================
\subsection{Condition de perpendicularité}
%====================================================

Le vecteur erreur :

\[
\vec{v} - \lambda \vec{u}
\]

doit être orthogonal à $\vec{u}$.

\bigskip

Donc :

\[
(\vec{v} - \lambda \vec{u}) \cdot \vec{u} = 0.
\]

---

%====================================================
\subsection{Démonstration de la formule}
%====================================================

On développe :

\[
\vec{v} \cdot \vec{u} - \lambda (\vec{u} \cdot \vec{u}) = 0.
\]

---

On obtient :

\[
\lambda = \frac{\vec{v} \cdot \vec{u}}{||\vec{u}||^2}.
\]

---

%====================================================
\subsection{Formule fondamentale}
%====================================================

\[
\boxed{
\text{proj}_{\vec{u}}(\vec{v}) =
\frac{\vec{v} \cdot \vec{u}}{||\vec{u}||^2} \vec{u}
}
\]

---

%====================================================
\subsection{Interprétation géométrique}
%====================================================

Le coefficient :

\[
\frac{\vec{v} \cdot \vec{u}}{||\vec{u}||}
\]

représente la longueur de la projection.

\bigskip

Ainsi :

\begin{center}
\textbf{projeter = prendre la composante de $\vec{v}$ dans la direction de $\vec{u}$}
\end{center}

---

%====================================================
\subsection{Application à la projection sur une droite}
%====================================================

On considère :

\begin{itemize}
\item une droite $(d)$ de vecteur directeur $\vec{u}$,
\item un point $A$,
\item un point $B$ sur la droite.
\end{itemize}

---

On a :

\[
\vec{BH} = \text{proj}_{\vec{u}}(\vec{BA})
\]

---

Donc :

\[
\boxed{
H = B + \frac{\vec{BA} \cdot \vec{u}}{||\vec{u}||^2} \vec{u}
}
\]

---

%====================================================
\subsection{Comparaison avec la méthode précédente}
%====================================================

\begin{itemize}
\item ancienne méthode : résoudre une équation
\item nouvelle méthode : formule directe
\end{itemize}

\bigskip

\begin{center}
\textbf{On gagne en rapidité et en clarté.}
\end{center}

---

%====================================================
\subsection{Cas particulier : vecteur unitaire}
%====================================================

Si $\vec{u}$ est unitaire ($||\vec{u}|| = 1$) :

\[
\boxed{
\text{proj}_{\vec{u}}(\vec{v}) = (\vec{v} \cdot \vec{u}) \vec{u}
}
\]

---

%====================================================
\subsection{Décomposition d’un vecteur}
%====================================================

Tout vecteur $\vec{v}$ peut s’écrire :

\[
\vec{v} =
\underbrace{\text{proj}_{\vec{u}}(\vec{v})}_{\text{partie parallèle}}
+
\underbrace{\left(\vec{v} - \text{proj}_{\vec{u}}(\vec{v})\right)}_{\text{partie orthogonale}}
\]

---

%====================================================
\subsection{Interprétation fondamentale}
%====================================================

\begin{center}
\textbf{Projection = décomposition d’un vecteur}
\end{center}

\bigskip

C’est exactement le principe qui sera utilisé en :

\begin{itemize}
\item algèbre linéaire,
\item mécanique (changement de base),
\item analyse vectorielle.
\end{itemize}

---

%====================================================
\subsection{Conclusion}
%====================================================

\begin{itemize}
\item la projection devient une opération vectorielle simple,
\item elle permet de décomposer un vecteur selon une direction,
\item elle prépare directement au changement de base.
\end{itemize}

\bigskip

\begin{center}
\textbf{On passe d’une méthode géométrique à une structure algébrique.}
\end{center}

%====================================================
\section{Lien avec le changement de base}
%====================================================

\subsection{Idée fondamentale}

Dans les parties précédentes, nous avons vu que projeter un vecteur consiste à extraire
sa composante selon une direction donnée.

\bigskip

Cette idée est en réalité au cœur du changement de base.

\bigskip

\begin{center}
\Large \textbf{Projeter un vecteur revient à lire une coordonnée dans une base adaptée.}
\end{center}

---

%====================================================
\subsection{Décomposition dans une base orthonormée}
%====================================================

Soit une base orthonormée $(\vec{e}_1,\vec{e}_2)$.

\bigskip

Tout vecteur $\vec{v}$ s’écrit :

\[
\vec{v} = v_1 \vec{e}_1 + v_2 \vec{e}_2
\]

---

\subsection{Calcul des coordonnées}

On prend le produit scalaire avec $\vec{e}_1$ :

\[
\vec{v} \cdot \vec{e}_1 = v_1
\]

car :

\[
\vec{e}_1 \cdot \vec{e}_1 = 1 \quad \text{et} \quad \vec{e}_2 \cdot \vec{e}_1 = 0
\]

---

\[
\boxed{
v_1 = \vec{v} \cdot \vec{e}_1
}
\]

---

%====================================================
\subsection{Interprétation}
%====================================================

On obtient :

\[
v_1 \vec{e}_1 = (\vec{v} \cdot \vec{e}_1)\vec{e}_1
\]

---

\begin{center}
\Large \textbf{C’est exactement la projection de $\vec{v}$ sur $\vec{e}_1$}
\end{center}

---

%====================================================
\subsection{Schéma}
%====================================================

\begin{center}
\begin{tikzpicture}[scale=1.2]

\draw[->] (0,0) -- (3,0) node[right] {$\vec{e}_1$};
\draw[->] (0,0) -- (0,2.5) node[above] {$\vec{e}_2$};

\draw[->] (0,0) -- (2,1.5) node[right] {$\vec{v}$};

\draw[dashed] (2,1.5) -- (2,0);

\node at (2,-0.4) {$v_1 \vec{e}_1$};

\end{tikzpicture}
\end{center}

---

%====================================================
\subsection{Lien avec la formule de projection}
%====================================================

On retrouve :

\[
\text{proj}_{\vec{e}_1}(\vec{v}) = (\vec{v} \cdot \vec{e}_1)\vec{e}_1
\]

---

\subsection{Cas général}

Si la base n’est pas orthonormée :

\[
\text{proj}_{\vec{u}}(\vec{v}) =
\frac{\vec{v} \cdot \vec{u}}{||\vec{u}||^2} \vec{u}
\]

---

%====================================================
\section{Interprétation matricielle}
%====================================================

\subsection{Vecteur comme colonne}

On écrit :

\[
\vec{v} =
\begin{pmatrix}
v_1 \\
v_2
\end{pmatrix}
\]

---

\subsection{Projection sur $\vec{e}_1$}

\[
\text{proj}_{\vec{e}_1}(\vec{v}) =
\begin{pmatrix}
v_1 \\
0
\end{pmatrix}
\]

---

\subsection{Matrice de projection}

\[
\boxed{
P =
\begin{pmatrix}
1 & 0 \\
0 & 0
\end{pmatrix}
}
\]

---

\[
P \vec{v} =
\begin{pmatrix}
v_1 \\
0
\end{pmatrix}
\]

---

%====================================================
\subsection{Interprétation}
%====================================================

\begin{center}
Projection = transformation linéaire
\end{center}

---

%====================================================
\section{Lien avec la mécanique}
%====================================================

Dans le chapitre sur le changement de base :

\[
\vec{v} = v_r \vec{e}_r + v_\theta \vec{e}_\theta
\]

---

On a :

\[
v_r = \vec{v} \cdot \vec{e}_r
\]

---

\begin{center}
\textbf{On projette la vitesse sur une direction}
\end{center}

---

\subsection{Conclusion physique}

\begin{itemize}
\item projeter = extraire une composante physique
\item changement de base = réécriture du vecteur
\end{itemize}

---

%====================================================
\section{Conclusion générale du chapitre}
%====================================================

\begin{center}
\Huge \textbf{Synthèse}
\end{center}

\bigskip

\begin{itemize}
\item la projection orthogonale permet de trouver une distance minimale
\item elle s’exprime via le produit scalaire
\item elle permet de décomposer un vecteur
\item elle correspond à une lecture de coordonnées
\item elle est liée aux matrices et au changement de base
\end{itemize}

\bigskip

\begin{center}
\Large \textbf{Projection = géométrie + algèbre}
\end{center}

\end{document}